\section{Limitations}

The evaluation covers only CHIP-8, PICO-8, and TIC-80. Three bundles provide
a heterogeneous design test but do not establish generality across all
virtual machines. PICO-8 and TIC-80 are bounded compatibility runtimes over
documented subsets; neither implementation claims exact upstream semantics.

Performance measurements, when sealed, will describe the tested hosts,
workloads, source revisions, package builds, and observation modes. They do
not establish universal performance, population-scale throughput, or
cross-host guarantees. The memory evidence is empirical for the tested
configurations and is not a proof of universally constant resident memory.
Persisted evidence can grow with run length even when reducer and recorder
buffers are bounded.

Python sink delivery is synchronous. This makes callback backpressure
explicit, but a slow consumer can affect execution time. Asynchronous or
buffered sink semantics are outside the current contract. Consumers that
retain values beyond a callback remain responsible for their own memory and
storage policies.

The work evaluates a machine-neutral execution and evidence interface; it
does not demonstrate substrate-independent computation. The common
orchestration contract does not imply common machine semantics, a complete
trace for every run, or a universal observation vocabulary.
